\section{Conclusion}
\label{sec:conclusion}

The discipline this paper asks for comes down to one refusal: not to read a corpus frequency as a population rate. A frequency is a score, and the distance between a score and the construct it stands for is the whole of the subject. What the corpus's own construction adds is construct-irrelevant variance that happens to run in a direction over time, which is why a before-and-after contrast can't be trusted on its own. Once that line holds the rest follows: a causal claim needs the target population, composition, measurement, and the units of variation and inference pulled apart, with parallel trends asked of the rate and of the filter both. Table~\ref{tab:reporting} is what that comes to on the page, and the ladder of Section~\ref{sec:diagnostics} is what it comes to in a decision.

The payoff isn't usually a clean effect. A corpus DiD is at its best when it pushes a causal claim downward, turning an apparent effect into a shared wave, a measurement shift, a composition artifact, or a verdict of not identified. Each is an answer worth reporting. The hard part is keeping that downgrade tied to evidence: a missed effect caused by thin data isn't a methodological success, which is why the failure map counts false negatives as carefully as false positives.

Before-and-after corpus data can carry a causal claim, but only under conditions worth stating plainly. The treatment has to be a named act at a named date; the comparison has to make a counterfactual plausible; composition and measurement have to be held fixed on pre-committed margins or shown not to move; and the result has to survive a sensitivity bound larger than the setting's plausible confounds. Where those hold, the claim is a bounded estimate scoped to the register measured, not a fact about the language at large.

The feminization case carried these points without resting on any of them. The same apparatus travels to other shocks and other languages, as far as the corpora reach. When a corpus frequency shifts after a shock, the discipline turns the question from whether the line moved to what moved: the rate, the filter, the measurement, or only the warrant for the word \mention{cause}.
